% !PRINTED 8
\weberpage{8}
If we therefore imagine $x_1, x_2, x_3, x_4$ placed in this order on the circumference of a circle, we always obtain a positive properly fractional $\varkappa^2$ when the $x_i$ follow one another in order of quantity in such a way that one begins with an arbitrary one as the smallest and then counts onward around the circle either to the right or to the left. We express this as follows:

\gesperrt{In order that $\varkappa^2$ be a positive proper fraction, the quantities $x_1, x_2, x_3, x_4$ must follow one another cyclically in order of quantity.}

§ 3. \gesperrt{Linear Transformation of the Elliptic Differential.}

If the two points $x_3$ and $x_4$ approach one and the same point $x$, then $x_3'$ and $x_4'$ approach the corresponding point $x'$. We then set $(x_3x_4):(x_3'x_4') = dx:dx'$ and obtain from (3), § 2, the relation between the differentials $dx, dx'$:
\weq{1}{\frac{(x_1x_2)\,dx}{(x_1x)(x_2x)} = \frac{(x_1'x_2')\,dx'}{(x_1'x')(x_2'x')}.}

In the same way, if in place of $x_1, x_2$ one puts two other points $x_3, x_4$, there results:
\weq{2}{\frac{(x_3x_4)\,dx}{(x_3x)(x_4x)} = \frac{(x_3'x_4')\,dx'}{(x_3'x')(x_4'x')},}

and if one multiplies and extracts the root:
\weq{3}{\frac{\sqrt{(x_1x_2)(x_3x_4)}\,dx}{\sqrt{(x_1x)(x_2x)(x_3x)(x_4x)}} = \frac{\sqrt{(x_1'x_2')(x_3'x_4')}\,dx'}{\sqrt{(x_1'x')(x_2'x')(x_3'x')(x_4'x')}},}

\gesperrt{hence a transformation of the elliptic differential by a linear substitution:}
\[
x = \frac{\alpha x' + \beta}{\gamma x' + \delta}.
\]

This substitution is completely determined when, for three arbitrary points $x_1, x_2, x_3$, the corresponding values $x_1', x_2', x_3'$ are arbitrarily given, and one obtains it from the equality of the cross ratio:
\weq{4}{\frac{(x_1x_2)(x_3x)}{(x_1x_3)(x_2x)} = \frac{(x_1'x_2')(x_3'x')}{(x_1'x_3')(x_2'x')},}

by solving for $x$; and the fourth value $x_4'$, the one belonging to $x_4$, is obtained from
%CONT>>
